% p157_1; source PDF page 157; manually verified editable table.
\begin{tikzpicture}[x=1cm,y=-1cm,inner sep=0pt,outer sep=0pt]
\path[use as bounding box] (0,0) rectangle (16.000,4.570);
\draw[line width=.4pt] (0.000,0.000) rectangle (2.800,0.855);
\node[align=center,text width=2.540cm,font={\linespread{1}\fontsize{10}{13.50}\selectfont\bfseries}] at (1.400,0.427) {角色};
\draw[line width=.4pt] (2.800,0.000) rectangle (16.000,0.855);
\node[align=center,text width=12.940cm,font={\linespread{1}\fontsize{10}{13.50}\selectfont\bfseries}] at (9.400,0.427) {要求};
\draw[line width=.4pt] (0.000,0.855) rectangle (2.800,1.710);
\node[align=left,text width=2.540cm,font={\linespread{1}\fontsize{10}{13.50}\selectfont}] at (1.400,1.282) {教唆犯};
\draw[line width=.4pt] (2.800,0.855) rectangle (16.000,1.710);
\node[align=left,text width=12.940cm,font={\linespread{1}\fontsize{10}{13.50}\selectfont}] at (9.400,1.282) {消除心理性贡献（打消犯意）};
\draw[line width=.4pt] (0.000,1.710) rectangle (2.800,2.565);
\node[align=left,text width=2.540cm,font={\linespread{1}\fontsize{10}{13.50}\selectfont}] at (1.400,2.138) {帮助犯};
\draw[line width=.4pt] (2.800,1.710) rectangle (16.000,2.565);
\node[align=left,text width=12.940cm,font={\linespread{1}\fontsize{10}{13.50}\selectfont}] at (9.400,2.138) {消除物理性、心理性的贡献};
\draw[line width=.4pt] (0.000,2.565) rectangle (2.800,4.570);
\node[align=left,text width=2.540cm,font={\linespread{1}\fontsize{10}{13.50}\selectfont}] at (1.400,3.567) {共同正犯};
\draw[line width=.4pt] (2.800,2.565) rectangle (16.000,4.570);
\node[align=left,text width=12.940cm,font={\linespread{1}\fontsize{10}{13.50}\selectfont}] at (9.400,3.567) {消除物理性、心理性的贡献。\\在预备阶段，要消除心理性贡献，需要明确告知退出意思。\\在实行阶段，要消除贡献，原则上要求有效防止既遂结果发生};
\end{tikzpicture}
